% Part: first-order-logic
% Chapter: sequent-calculus
% Section: quantifier-rules

\documentclass[../../../include/open-logic-section]{subfiles}

\begin{document}

\olfileid[pt]{fol}{seq}{qrl}

\olsection{Regras de Quantificadores}

\subsection{Regras para $\lforall$}

\begin{defish}
\Axiom$ !A(t), \Gamma \fCenter \Delta$
\RightLabel{\LeftR{\lforall}}
\UnaryInf$ \lforall[x][!A(x)],\Gamma \fCenter \Delta$
\DisplayProof
\hfill
\Axiom$ \Gamma \fCenter \Delta, !A(a) $
\RightLabel{\RightR{\lforall}}
\UnaryInf$ \Gamma \fCenter \Delta, \lforall[x][!A(x)]$
\DisplayProof
\end{defish}

Em \LeftR{\lforall}, $t$ é um termo fechado (isto é, um termo sem
variáveis). Em \RightR{\lforall}, $a$~é !!a{constant} que não pode
ocorrer em lugar algum do sequente inferior da regra \RightR{\lforall}.
Chamamos $a$ de \emph{variável própria} da inferência
\RightR{\forall}.\footnote{Usamos o termo ``variável própria'' embora $a$
na regra acima seja !!a{constant}. Isso tem razões históricas.}

\subsection{Regras para $\lexists$}

\begin{defish}
\Axiom$ !A(a), \Gamma \fCenter \Delta $
\RightLabel{\LeftR{\lexists}}
\UnaryInf$ \lexists[x][!A(x)], \Gamma \fCenter \Delta$
\DisplayProof
\hfill
\Axiom$ \Gamma \fCenter \Delta, !A(t) $
\RightLabel{\RightR{\lexists}}
\UnaryInf$ \Gamma \fCenter \Delta, \lexists[x][!A(x)]$
\DisplayProof
\end{defish}

Novamente, $t$~é um termo fechado, e $a$~é !!a{constant} que não
ocorre no sequente inferior da regra \LeftR{\lexists}. Chamamos $a$ de
\emph{variável própria} da inferência \LeftR{\lexists}.

A condição de que uma variável própria não ocorra no sequente inferior
da inferência \RightR{\lforall} ou \LeftR{\lexists} é chamada de
\emph{condição de variável própria}.

\begin{explain}
Lembre-se da convenção de que, quando $!A$ é !!a{formula} com a
!!{variable}~$x$ livre, indicamos isso escrevendo~$!A(x)$. No
mesmo contexto, $!A(t)$ é então uma abreviação de~$\Subst{!A}{t}{x}$. Assim,
poderíamos também escrever a regra \RightR{\lexists} como:
\begin{prooftree}
  \Axiom$\Gamma \fCenter \Delta, \Subst{!A}{t}{x}$
  \RightLabel{\RightR{\lexists}}
  \UnaryInf$\Gamma \fCenter \Delta, \lexists[x][!A]$
\end{prooftree}
Note que $t$ pode já ocorrer em~$!A$, por exemplo, $!A$~poderia
ser~$\Atom{\Obj P}{t,x}$. Assim, inferir $\Gamma \Sequent \Delta,
\lexists[x][\Atom{\Obj P}{t,x}]$ de~$\Gamma \Sequent \Delta,
\Atom{\Obj P}{t,t}$ é uma aplicação correta
de~\RightR{\lexists}---pode-se ``substituir'' uma ou mais, e não
necessariamente todas, as ocorrências de~$t$ na premissa pela
!!{variable}~$x$ ligada. Contudo, as condições de variável própria em
\RightR{\lforall} e~\LeftR{\lexists} exigem que a
!!{constant}~$a$ não ocorra em~$!A$. Portanto, não se pode
inferir corretamente $\Gamma \Sequent \Delta, \lforall[x][\Atom{\Obj P}{a,x}]$ de
$\Gamma \Sequent \Delta, \Atom{\Obj P}{a,a}$ usando~$\RightR{\lforall}$.
\end{explain}

\begin{explain}
Em \RightR{\lexists} e \LeftR{\lforall} não há restrições sobre
o termo~$t$. Por outro lado, nas regras \LeftR{\lexists} e
\RightR{\lforall}, a condição de variável própria exige que a
!!{constant}~$a$ não ocorra em lugar algum fora de~$!A(a)$ no
sequente superior. Isso é necessário para garantir que o sistema seja correto,
isto é, que só !!{derive} sequentes que sejam válidos. Sem essa
condição, o seguinte seria permitido:
\begin{prooftree}
  \Axiom$!A(a) \fCenter !A(a)$
  \RightLabel{*\LeftR{\lexists}}
  \UnaryInf$\lexists[x][!A(x)] \fCenter !A(a)$
  \RightLabel{\RightR{\lforall}}
  \UnaryInf$\lexists[x][!A(x)] \fCenter \lforall[x][!A(x)]$
  \DisplayProof\bottomAlignProof
  \qquad
  \Axiom$!A(a) \fCenter !A(a)$
  \RightLabel{*\RightR{\lforall}}
  \UnaryInf$!A(a) \fCenter \lforall[x][!A(x)]$
  \RightLabel{\LeftR{\lexists}}
  \UnaryInf$\lexists[x][!A(x)] \fCenter \lforall[x][!A(x)]$
\end{prooftree}
Contudo, $\lexists[x][!A(x)] \Sequent \lforall[x][!A(x)]$ não é válido.
\end{explain}

\end{document}
